\subsection{从李代数到李群}

命题 1. 设 G 是单位元为 e 的李群芽。G 中存在一个由 G 的李子群组成的 e 的开邻域基。

给 \(\mathrm{L}(\mathrm{G})\) 赋予一个与其拓扑相容的范数，并使得对 \(\| [x,y]\| \leqslant \| x\| \| y\|\) 中所有 \(x,y\) 都有 \(\mathrm{L}(\mathrm{G})\)。令 \(\mathrm{G}_1\) 为由 \(\mathrm{L}(\mathrm{G})\) 定义的李群。由 § 4，第 2 号，定理 2，G 与 \(\mathrm{G}_1\) 局部同构。因此只需应用 § 4，第 2 号，引理 3 (iii)。

定理 1. 设 \(\mathbf{L}\) 是完备可赋范李代数。存在一个李群 \(\mathbf{G}\)，使得 \(\mathbf{L}(\mathbf{G})\) 同构于 \(\mathbf{L}\)。任意两个这样的群都局部同构。

第一个断言已在 § 4，第 2 号，引理 3 中证明。第二个断言是 § 4，第 2 号，定理 2 的特例。

定理 2. 设 G 是李群，b 是 L(G) 的具有拓扑补空间的李子代数。存在 G 的李子群 H，使得 L(H) = b。若 H₁ 和 H₂ 是 G 的李子群，且 L(H₁) = L(H₂) = b，则 H₁ ∩ H₂ 在 H₁ 和 H₂ 中都是开集。

第一个断言由命题 1 和 §4，第 2 号，定理 3 得到。第二个断言是 §4，第 2 号，定理 3 的特例。

定理 3. 设 \(\mathbf{G}\) 和 \(\mathbf{H}\) 是李群，\(h\) 是从 \(\mathbf{L}(\mathbf{G})\) 到 \(\mathbf{L}(\mathbf{H})\) 的连续态射。

\begin{enumerate}
\def\labelenumi{(\roman{enumi})}
\item
  存在 \(\mathbf{G}'\) 的开子群 \(\mathbf{G}\) 以及从 \(\phi\) 到 \(\mathbf{G}'\) 的李群态射 \(\mathbf{H}\)，使得 \(h = \mathbf{L}(\phi)\)。
\item
  设 \(\mathbf{G}_1, \mathbf{G}_2\) 是 \(\mathbf{G}\) 的开子群，\(\phi_i\) 是从 \(\mathbf{G}_i\) 到 \(\mathbf{H}\) 的态射，且 \(h = \mathbf{L}(\phi_i)\)。则 \(\phi_1\) 与 \(\phi_2\) 在 \(\mathbf{G}\) 的某个开子群上相同。
\end{enumerate}

由命题 1，这由 § 4，第 1 号，定理 1 得到。

命题 2. 设 \(\mathbf{G}\) 是李群，\(\mathfrak{h}\) 是 \(\mathbf{L}(\mathbf{G})\) 的具有拓扑补空间的李子代数。下列条件等价：

\begin{enumerate}
\def\labelenumi{(\roman{enumi})}
\item
  存在 \(\mathbf{G}'\) 的开子群 \(\mathbf{G}\) 以及 \(\mathbf{H}\) 的正规李子群 \(\mathbf{G}'\)，使得 \(\mathrm{L(H)} = \mathfrak{h}\)。
\item
  \(\mathfrak{h}\) 是 \(\mathrm{L}(\mathrm{G})\) 的理想。
\end{enumerate}

若存在具有 (i) 中性质的 \(G'\)\(H\) 和 H，则由 § 3，第 12 号，命题 47，有 \(L(G') = L(G)\)，且 \(L(H)\) 是 \(L(G')\) 的理想。

假设 \(\mathfrak{h}\) 是 \(\mathrm{L}(\mathrm{G})\) 的理想。存在一个李群 \(\mathrm{F}\)，使得 \(\mathrm{L}(\mathrm{F}) = \mathrm{L}(\mathrm{G}) / \mathfrak{h}\)（定理 1）。令 h 为从 \(h\)\(\mathrm{L}(\mathrm{G})\) 到 \(\mathrm{L}(\mathrm{F})\) 的典范态射。由定理 3 (i)，存在 \(\mathrm{G}'\) 的开子群 \(\mathrm{G}\) 以及从 \(\phi\) 到 \(\mathrm{G}'\) 的李群态射 \(\mathrm{F}\)，使得 \(\mathrm{L}(\phi) = \mathrm{h}\)。由 § 3，第 8 号，\(\mathrm{H}\) 的核 \(\phi\) 是 \(\mathrm{G}'\) 的李子群，且 \(\mathrm{L}(\mathrm{H}) = \mathrm{Ker} \, \mathrm{L}(\phi) = \mathrm{Ker} \, h = \mathfrak{h}\)。最后，由于 ，\(\mathrm{H}\) 是 \(\mathrm{G}'\)\(\mathrm{H} = \mathrm{Ker} \, \phi\) 的正规子群。

